%%%%

%%%   Самарский государственный технический университет

%%%   Кафедра прикладной математики и информатики

%%%

%%%   Преамбула для статьи в журнал "Вестник Самарского государственного

%%%   технического университета. Серия: «Физико-математические науки»"

%%%   Ориентирована под MikTEx 2.4 for Win32

%%%

%%%   Хотя сам журнал собирается под GNU/Linux на дистрибутиве TeXLive



\documentclass[11pt,a4paper,twoside]{article}



\usepackage[cp1251]{inputenc}

\usepackage[T2A]{fontenc}

\usepackage[english,russian]{babel}

\usepackage{samgtubib}

\usepackage[dvips]{graphicx}

\usepackage[multiple]{footmisc}



\usepackage{samgtu}

\renewcommand{\VestnikYear}{2009} % Год издания Вестника

\renewcommand{\VestnikNumber}{2\,(19)} % Номер Вестника

\renewcommand{\VestnikMonth}{Ноябрь} % Месяц выхода Вестника

\renewcommand{\VestnikData}{01.11.09} % Дата подписания Вестника в печать

\renewcommand{\VestnikUPL}{26,64} % Усл. печ. л.

\renewcommand{\VestnikUIL}{25,73} % Уч.-изд. л.

\renewcommand{\VestnikRegNumber}{479/09} % Регистрационный номер

\renewcommand{\VestnikOrder}{994} % Номер заказа







%

%\begin{document}

%

%\Partition{Алгебра}{}



\renewcommand{\baselinestretch}{0.86}



\renewcommand{\firstpage}{7}

\renewcommand{\lastpage}{14}



\udk{512.572}

\msc{16R10, 17A32, 17B01, 17B63}

\title{Многообразия линейных алгебр полиномиального~роста}{Многообразия линейных алгебр полиномиального роста}



\author{О.~И.~Череватенко}% Автор(ы) для заголовка, если авторов несколько укать \inst

{О.~И.~Ч\,е\,р\,е\,в\,а\,т\,е\,н\,к\,о}% Автор(ы) для оглавления и колонтитула через \,



\institute{\ulgpu.}



\email{E-mail}{chai@pisem.net}



\otherinf{\fullauthor{Ольга Ивановна Череватенко} \regalia{к.ф.-м.н.} \position{доцент} \dept{каф. высшей математики} }



\keywords{ассоциативная алгебра, алгебра Пуассона, алгебра Ли, многообразие алгебр, рост многообразия}



\workabstract{

Приведён обзор результатов о многообразиях линейных алгебр полиномиального роста. 

Приводятся эквивалентные условия полиномиальности роста многообразий ассоциативных алгебр, многообразий алгебр Ли, многообразий алгебр Лейбница, многообразий алгебр Пуассона и многообразий алгебр Лейбница"--~Пуассона. 

Показано, что при изучении многообразий линейных алгебр полиномиального роста важную роль играют многообразия почти полиномиального роста.

}



\engtitle{Varieties of linear algebras of polynomial growth}



\engauthor{O.~I.~Cherevatenko}{O.~I.~C\,h\,e\,r\,e\,v\,a\,t\,e\,n\,k\,o}





\enginstitute{\engulgpu.}



\otherenginf{\fullauthor{Olga I. Cherevatenko} \regalia{Ph.\,D. Phys. \& Math.} \position{Associate Professor} \dept{Dept. of Higher Mathematics}}





\engabstract{The paper is survey of results of investigations on varieties of linear algebras of polynomial growth. We give equivalent conditions of the polynomial codimension growth of a variety of associative algebras, Lie algebras, Leibniz algebras, Poisson algebras, Leibniz-Poisson algebras. It is shown that in the study of varieties of linear algebras of polynomial growth varieties of almost polynomial growth play an important role.

}





\engkeywords{associative algebra, Poisson algebra, Lie algebra, variety of algebras, growth of a variety}



\date{27/IX/2013} % Дата поступления статьи в редакцию.

\revisiondate{12/X/2013} % В окончательном виде



\maketitle





{\bf Введение.} На протяжении всей работы предполагается, если это специально не оговорено, что основное поле имеет нулевую характеристику.



Пусть $K(X)$ --- свободная алгебра над полем $K$, где

$X=\{x_1,x_2, \dots\}$ "--- счётное множество свободных образующих.

Обозначим через $A$ некоторую $PI$-алгебру. Совокупность всех

тождеств алгебры $A$ образует $T$-идеал $Id(A)$ в свободной

алгебре $K(X)$. Пусть ${\bf V}$ "--- многообразие алгебр, порождённое

алгеброй $A$. Тогда алгебра $K(X,{\bf V})=K(X)/Id(A)$ является

относительно свободной алгеброй многообразия ${\bf V}$.



Пусть $P_n$ "--- подпространство в пространстве $K(X)$, состоящее

из всех полилинейных элементов степени $n$ от переменных

$x_1$, $x_2$, $\dots,x_n$. В случае основного поля нулевой характеристики вся

информация о многообразии ${\bf V}$ содержится в его полилинейных

компонентах $P_n({\bf V})=P_n/(P_n \cap Id({\bf V}))$, $n=1, 2, \dots$.

Асимптотическое поведение последовательности $c_n({\bf V}){=}\mathop{\rm dim} P_n({\bf V}),$

$n=1, 2, \dots,$ называют ростом многообразия ${\bf V}$. 

Говорят, что многообразие ${\bf V}$ имеет полиномиальный рост, если существуют такие константы $C$ и $k$, что для любого $n$ выполнено неравенство $c_n({\bf V})\le Cn^k$.



Пространство $P_n({\bf V})$ наделено структурой левого $S_n$-модуля,

 где $S_n$ "--- симметрическая группа степени

 $n$. Напомним, что последовательность $\lambda=

(\lambda_1,\lambda_2, \dots ,\lambda_k)$ называют {\it разбиением

числа $n$} и обозначают $\lambda \vdash n,$ если $\lambda_1 +

\lambda_2+  \dots  + \lambda_k=n$ и $\lambda_1\ge \lambda_2 \ge

 \dots  \ge \lambda_k>0.$ Пусть $\chi_{\lambda}$ "--- характер неприводимого представления

симметрической группы, соответствующий разбиению $\lambda$ числа

$n$. Тогда в силу вполне приводимости модуля $P_n({\bf V})$ для многообразия ${\bf V}$ имеет место

разложение

\begin{equation}

\label{dop10}

   \chi_n({\bf V})=\chi_n(P_n({\bf V})) = \sum_{\lambda \vdash n}m_{\lambda}({\bf V}) \chi_{\lambda},

\end{equation}

где $m_{\lambda}({\bf V})$ "--- степени неприводимых представлений,

соответствующих разбиению $\lambda$ числа $n$.

{\it Кодлина} многообразия определяется как сумма

$$

l_n({\bf V})=\sum_{\lambda \vdash n}m_{\lambda}({\bf V}).

$$

Если $\lambda$ "--- разбиение некоторого числа $n$, то будем

обозначать символом $\lambda' $ сопряженное разбиение к

$\lambda$.



В следующей теореме приводится критерий полиномиального роста для общего случая.



\begin{theorem}[1 \cite{cherev:bib:G-MSP}]

Пусть ${\bf V}$ "--- некоторое многообразие линейных алгебр. Последовательность $\{c_n({\bf V})\}_{n \geq 1}$ ограничена полиномом тогда и только тогда$,$ когда выполнены следующие условия\/{\rm :}

\begin{itemize}

\item[\rm (i)] существует такая константа $C,$ что в сумме \eqref{dop10}

 $m_\lambda({\bf V}) = 0$ в случае$,$ если либо выполнено условие $n -

\lambda_1 > C,$ либо $n -\lambda' _1 > C;$



\item[\rm (ii)] существуют такие константы $C$ и $k,$ что для любого $n$ существует $m\leq Cn^k$ расстановок скобок $T_1,$ $T_2,$ $\dots ,$ $T_m$ таких$,$ что любой элемент $f\in P_n({\bf V})$ может быть записан как

$$

f = \sum_{i=1}^mf_i \quad (\mathop{mod} Id({\bf V})),

$$

где $f_i \in P_n^{T_i}({\bf V}).$



\end{itemize}

\end{theorem}





Далее будут приведены более конкретные эквивалентные условия полиномиального роста для конкретных случаев.



Договоримся опускать скобки при их левонормированной расстановке: $$((ab)c) = abc.$$



%====================================================================================

%====================================================================================

%====================================================================================

\smallskip



{\bf Ассоциативные алгебры.} В теории ассоциативных алгебр очень важную роль играет бесконечно порожденная алгебра Грассмана $\Lambda$ и алгебра верхнетреугольных матриц порядка 2, которую обозначим через $UT_2$.



\vskip 0.1in



\begin{theorem}[2 \cite{cherev:bib:Dr}]

Для алгебры $\Lambda$ верны следующие утверждения\/{\rm:}



\begin{itemize}

\item[\rm (i)] полилинейное тождество $ [x,y,z]=0$

порождает идеал тождеств алгебры Грассмана $\Lambda;$



\item[\rm (ii)] рост многообразия $\mathop{var}(\Lambda),$ порождённого алгеброй $\Lambda,$

является почти полиномиальным, причём $c_n(\Lambda)=2^{n-1};$



\item[\rm (iii)] базис полилинейной компоненты $P_n(\Lambda)$ состоит из элементов вида

$$

 [x_{i_1},x_{i_2}] \dots 

[x_{i_{2p-1}},x_{i_{2p}}] x_{j_1} x_{j_2} \dots  x_{j_q},

$$

где $\{i_1, \dots ,i_{2p},j_1, \dots ,j_q\}=\{1, \dots ,n\},$

$i_1<i_2< \dots <i_{2p},$ $j_1< \dots <j_q.$



\end{itemize}



\end{theorem}



\smallskip



\begin{theorem}[3  \cite{cherev:bib:Dr}]

Для алгебры $UT_2$ верны следующие утверждения\/{\rm:}

\begin{itemize}





\item[\rm (i)] полилинейное тождество $[x_1,x_2] [x_3,x_4]=0$

порождает идеал тождеств алгебры $UT_2;$



\item[\rm (ii)] рост многообразия $var(UT_2),$ порождённого алгеброй $UT_2,$

является почти полиномиальным$,$ причём $c_n(UT_2)=2^{n-1}(n-2)+2;$



\item[\rm(iii)] базис полилинейной компоненты $P_n(UT_2)$ состоит из элементов вида

$$

 [x_{i_1},x_{i_2}, \dots ,x_{i_p}] x_{j_1} x_{j_2}  \dots  x_{j_q},

$$

 где

$\{i_1, \dots ,i_p,j_1, \dots ,j_q\}=\{1, \dots ,n\}$, $j_1< \dots <j_q$ и

переменные в коммутаторе $[x_{i_1}, \dots ,x_{i_p}]$ упорядочены

следующим образом\/{\rm :}

$$i_1> i_2<i_3< \dots <i_p.$$



\end{itemize}



\end{theorem}



\smallskip



\begin{theorem}[4 \cite{cherev:bib:Kemer}] Для многообразия ассоциативных алгебр ${\bf V}$ следующие

условия эквивалентны\/{\rm:}

\begin{itemize}

\item[\rm (i)] последовательность $\{c_n({\bf V})\}_{n \geq 1}$ ограничена

полиномом\/{\rm ;}



\item[\rm (ii)] $UT_2 \notin {\bf V},$ $\Lambda \notin {\bf V};$



\item[\rm (iii)]  существует такая константа $C,$ что в сумме \eqref{dop10}  $m_\lambda({\bf V}) = 0$ в случае$,$ если выполнено условие $n - \lambda_1 > C;$



\item[\rm (iv)] кодлина многообразия ${\bf V}$ является конечной$.$



\end{itemize}



\end{theorem}



%====================================================================================

%====================================================================================

%====================================================================================



\smallskip



{\bf Алгебры Ли.} Пусть ${\bf N}_s{\bf A}$ "--- многообразие алгебр Ли, определяемое тождеством

$[[x_1,x_2], \dots ,[x_{2s+1},x_{2s+2}]]=0.$

В работе \cite{cherev:bib:M2} показано, что многообразие ${\bf N}_2{\bf A}$ имеет почти полиномиальный рост.



\smallskip



\begin{theorem}[5 \cite{cherev:bib:B-Z, cherev:bib:M5}]

Для многообразия алгебр Ли ${\bf V}$ следующие условия эквивалентны\/$:$





\begin{itemize}



\item[\rm (i)] многообразие ${\bf V}$ имеет полиномиальный рост$;$



\item[\rm (ii)] для некоторого $s$ выполнено условие ${\bf N}_2{\bf A} \not\subseteq {\bf V} \subset {\bf N}_s{\bf A};$



\item[\rm (iii)]  существует такая константа $C,$ что в сумме \eqref{dop10}  $m_\lambda({\bf V}) = 0$ в случае$,$ если выполнено условие $n - \lambda_1 > C.$





\end{itemize}



\end{theorem}



\smallskip



Обозначим через ${\bf A}^2$ метабелево многообразие алгебр Ли, которое определяется тождеством

$[[x_1,x_2],[x_3,x_4]]=0.$

Хорошо известно, что многообразие ${\bf A}^2$ является наименьшим многообразием

среди всех многообразий алгебр Ли роста не ниже

полиномиального, то есть если рост некоторого многообразия ${\bf V}$

не ниже полиномиального, то ${\bf A}^2 \subseteq {\bf V}$.





%====================================================================================

%====================================================================================

%====================================================================================

\smallskip



{\bf Алгебры Лейбница.} Алгебры Лейбница "--- неантикоммутативные алгебры Ли, которые определяются тождеством Лейбница

$$

[[x,y],z]=[[x,z],y]+[x,[y,z]].

$$

Тождество Лейбница представляет собой следующее условие: оператор правого умножения на любой элемент алгебры является дифференцированием.



Обозначим через $\widetilde{{\bf N}_s {\bf A}}$ многообразие алгебр Лейбница, определяемое тождеством

$$

   [[x_1,x_2], \dots ,[x_{2s+1},x_{2s+2}]]=0.

$$

Понятно, что многообразие $\widetilde{{\bf N}_s {\bf A}}$ с тождеством $[x,x]=0$ "--- в точности многообразие ${\bf N}_s{\bf A}$.



Обозначим через $\widetilde{\bf V}_1$ многообразие алгебр Лейбница, определенное тождеством $[x_1,[x_2,x_3],[x_4,x_5]]=0.$

Данное многообразие имеет почти полиномиальный рост (см. \cite{cherev:bib:M-V}).



\smallskip



\begin{theorem}[6 \cite{cherev:bib:M-Ch, cherev:bib:M-Ch2}] 

Для многообразия алгебр Лейбница ${\bf V}$ следующие условия эквивалентны\/$:$



\begin{itemize}



\item[\rm (i)] многообразие ${\bf V}$ имеет полиномиальный рост\/$;$



\item[\rm (ii)] для некоторого $s$ выполнено условие 

$$

{\bf N}_2{\bf A},\widetilde{\bf V}_1\not\subseteq {\bf V} \subset

\widetilde{{\bf N}_s {\bf A}};

$$



\item[\rm (iii)]  существует такая константа $C,$ что в сумме \eqref{dop10}  $m_\lambda({\bf V}) = 0$ в случае$,$ если выполнено условие $n - \lambda_1 > C.$





\end{itemize}



\end{theorem}



\smallskip



Пусть ${\bf B}$ "--- многообразие алгебр Лейбница, определённое тождеством

$$

[x, [y,z]]=0.

$$

В работе \cite{cherev:bib:ChOI} показано, что многообразия ${\bf A}^2$ и ${\bf B}$ исчерпывают весь список минимальных многообразий среди всех многообразий алгебр Лейбница роста не ниже полиномиального, то есть если рост некоторого многообразия ${\bf V}$ не ниже полиномиального, то либо ${\bf A}^2 \subseteq {\bf V}$, либо ${\bf B} \subseteq {\bf V}$.



%====================================================================================

%====================================================================================

%====================================================================================

\smallskip



{\bf Алгебры Пуассона.} Алгебра $A=A({}+{},{}\cdot{},\{ ~ , ~ \}, K)$ над произвольным полем $K$

называется алгеброй Пуассона, если $A({}+{}, {}\cdot{} , K)$ "--- ассоциативная коммутативная алгебра с~единицей, $A(+,\{ ~ , ~ \},K)$ "--- алгебра Ли с операцией умножения \mbox{\{{ },{ }\}}, которая называется скобкой Пуассона, и выполняется правило Лейбница:

$$

\{a \cdot b,c\}=a\cdot \{b,c\} + \{ a, c\} \cdot b, \qquad a,b,c

\in A.

$$





Пусть $\Lambda$ "--- бесконечномерная алгебра Грассмана с единицей и операцией умножения $\wedge$. Введём в алгебре $\Lambda$ два новых умножения:

\begin{equation}

\label{poisson_eq10}

a\cdot b = \frac{1}{2}(a\wedge b+b\wedge a), \quad 

\{a,b\}=a\wedge b-b\wedge a, \qquad a, b \in \Lambda.

\end{equation}

Нетрудно проверить, что алгебра $(\Lambda,+,\cdot, \{ ~ , ~\})$ будет алгеброй Пуассона, которую обозначим через $G$.



\smallskip





\begin{theorem}[7  \cite{cherev:bib:MPR}]

Для алгебры $G$ верны следующие утверждения\/$:$



\begin{itemize}



\item[\rm (i)] полилинейное тождество $\{x,y,z\}=0$  порождает идеал тождеств алгебры Пуассона $G;$



\item[\rm (ii)] рост многообразия $\mathop{var}(G),$ порождённого алгеброй $G,$ является почти полиномиальным$,$ причём $c_n(G)=2^{n-1};$



\item[\rm (iii)] базис полилинейной компоненты $P_n(G)$ состоит из элементов вида

$$

 \{x_{i_1},x_{i_2}\}\cdot \dots \cdot

\{x_{i_{2p-1}},x_{i_{2p}}\}\cdot x_{j_1}\cdot x_{j_2}

\cdot \dots \cdot x_{j_q},

$$

где $\{i_1, \dots ,i_{2p},j_1, \dots ,j_q\}=\{1, \dots ,n\},$

$i_1<i_2< \dots <i_{2p},$ $j_1< \dots <j_q.$





\end{itemize}



\end{theorem}



\smallskip



Пусть $A_L$ "--- некоторая ненулевая алгебра Ли с лиевым умножением $[~ , ~]$ над произвольным полем $K$. 

Рассмотрим векторное пространство $A = A_L \oplus K,$ в котором определим операции $\cdot$ и $\{~,~\}$ следующим образом: 

\begin{equation}

\label{rps10}

(a+\alpha) \cdot (b+\beta) = (\beta a + \alpha b) + \alpha \beta, \atop

\{a+\alpha, b+\beta \} = [a,b], \quad a,b \in A_L, \quad \alpha,\beta \in K.

\end{equation}

Тогда полученная алгебра $(A, {}+{}, {}\cdot{} ,\{~,~\}, K)$ будет являться алгеброй Пуассона.



Пусть $U^L_2$ "--- двумерная метабелева алгебра Ли с базисом $\{a, b\}$ и таблицей умножения $[a,b] = a$, $[b,a]=-a$, $[a,a]=[b,b]=0$. 

Обозначим через $U_2$  алгебру Пуассона $U^L_2 \oplus K$ с операциями \eqref{rps10}.







\smallskip



\begin{theorem}[8 \cite{cherev:bib:RSM12}]

Для алгебры $U_2$ верны следующие утверждения\/$:$



\begin{itemize}

\item[\rm (i)] полилинейные тождества

$$

 \{x_1,x_2\} \cdot \{x_3,x_4\}=0, \quad 

\{\{x_1,x_2\}, \{x_3,x_4\}\}=0

$$

порождают идеал тождеств алгебры Пуассона $U_2;$



\item[\rm (ii)] рост многообразия $var(U_2),$ порождённого алгеброй $U_2,$

является почти полиномиальным$,$ причём $c_n(U_2)=2^{n-1}(n-2)+2;$



\item[\rm (iii)] базис полилинейной компоненты $P_n(U_2)$ состоит из элементов вида

$$

 \{x_{i_1},x_{i_2}, \dots ,x_{i_p}\}\cdot x_{j_1}\cdot x_{j_2} \cdot \dots \cdot x_{j_q},

$$

 где

$\{i_1, \dots ,i_p,j_1, \dots ,j_q\}=\{1, \dots ,n\}$, $j_1< \dots <j_q$ и

переменные в мономе $\{x_{i_1}, \dots ,x_{i_p}\}$ упорядочены следующим образом\/$:$

$$

i_1> i_2<i_3< \dots <i_p.

$$





\end{itemize}



\end{theorem}



\smallskip



Следующая теорема показывает, что в случае алгебр Пуассона алгебры $U_2$ и $G$ играют ту же роль, что и~алгебры $UT_2$ и $\Lambda$ в~ассоциативном случае.



\smallskip



\begin{theorem}[9 \cite{cherev:bib:RSM12}]

Для многообразия алгебр Пуассона ${\bf V}$ следующие условия эквивалентны\/$:$



\begin{itemize}

\item[\rm (i)] последовательность $\{c_n({\bf V})\}_{n \geq 1}$ ограничена полиномом$;$



\item[\rm (ii)] $U_2 \notin {\bf V},$ $G \notin {\bf V};$



\item[\rm (iii)]  существует такая константа $C,$ что в сумме \eqref{dop10}  $m_\lambda({\bf V}) = 0$ в случае$,$ если выполнено условие $n - \lambda_1 > C;$



\item[\rm (iv)] кодлина многообразия ${\bf V}$ является конечной$.$



\end{itemize}





\end{theorem}



\smallskip



Пусть $N^L_3$ "--- трёхмерная нильпотентная алгебра Ли с базисом $\{a,b,c\}$ и таблицей умножения  $[b,a] = c$, $[a,c] = [b,c] = 0$. Обозначим через $N_3$  алгебру Пуассона $N^L_3 \oplus K$ с операциями \eqref{rps10}.



Пусть $\Lambda_{2n}$ "--- алгебра Грассмана с единицей и $2n$ образующими элементами $\{e_1, \dots ,e_{2n}\}$. Определим в алгебре $\Lambda_{2n}$ операции умножения (\ref{poisson_eq10}).

Полученную алгебру Пуассона обозначим через $G_{2n}$.



Следующая теорема показывает, что многообразие, порождённое либо алгеброй $N_3$, либо алгеброй $G_2$, является наименьшим многообразием алгебр Пуассона в классе всех многообразий алгебр Пуассона, имеющих рост не ниже полиномиального.





\smallskip



\begin{theorem}[10 \cite{cherev:bib:RSM13}]

Для алгебр $N_3$ и $G_2$ верны следующие утверждения\/$:$



\begin{itemize}



\item[\rm (i)] $\mathop{var}(N_3) = \mathop{var}(G_2) = {\bf V};$



\item[\rm (ii)] тождества $\{x_1,x_2,x_3\}=0,$ $\{x_1,x_2\}\cdot \{x_3,x_4\}=0$

порождают идеал тождеств многообразия ${\bf V};$



\item[\rm (iii)] для любого $n$ выполнено равенство $c_n({\bf V}) = 1+{n(n-1)}/{2};$



\item[\rm (iv)] многообразие ${\bf V}$ является наименьшим многообразием среди всех многообразий алгебр Пуассона роста не ниже полиномиального$,$ то есть если рост некоторого многообразия ${\bf W}$

не ниже полиномиального$,$ то ${\bf V} \subseteq {\bf W}.$



\end{itemize}



\end{theorem}



\smallskip



{\bf Алгебры Лейбница"--~Пуассона.} В данном пункте рассматриваются алгебры Пуассона с~неантикоммутативной операцией $\{~,~\}$, которые будем называть алгебрами Лейбница"--~Пуассона. Более точно векторное пространство $A$ над полем $K$ с двумя $K$-биллинейными операциями умножения ${}\cdot{}$ и $\{~,~\}$ называется алгеброй Лейбница"--~Пуассона, если относительно операции ${}\cdot{}$ пространство $A$ является коммутативной ассоциативной алгеброй с единицей, относительно операции $\{~,~\}$ "--- алгеброй Лейбница, и данные операции связаны правилами

$$

\{a \cdot b,c\}=a\cdot \{b,c\} +  \{ a, c\} \cdot b,

\quad 

\{c, a \cdot b\}=a\cdot \{c,b\} +  \{ c,a\} \cdot b,

$$

где $a,b,c\in A$.



Обозначим через $\Gamma_n$ пространство в свободной алгебре Лейбница"--~Пуассона, состоящее

из всех полилинейных элементов степени $n$ от переменных

$x_1, \dots ,x_n$ и являющееся линейной оболочкой элементов вида

$$

 \{x_{i_1}, \dots ,x_{i_s}\}\cdot \dots \cdot

\{x_{j_1}, \dots ,x_{j_t}\}, ~~~s\geq 2, \dots ,~t \geq 2.

$$



\smallskip



\begin{theorem}[11 \cite{cherev:bib:RSM14}]

Пусть ${\bf V}$ "--- нетривиальное многообразие алгебр Лейбница"--~Пуассона над произвольным полем$.$ Тогда либо



\begin{itemize}



\item[\rm (i)] $c_n({\bf V}) \geq 2^{n-1}$ для любого $n,$



\noindent либо





\item[\rm (ii)] найдётся такой многочлен $f(x)\in \mathbb{Q}[x],$ что для всех достаточно больших $n$ будет выполнено равенство $c_n({\bf V}) = f(n).$



\end{itemize}



\end{theorem}



\smallskip



Пусть ${\bf V}$ "--- некоторое многообразие алгебр Лейбница"--~Пуассона$,$ $\mathop{Id}({\bf V})$ "--- идеал тождеств многообразия ${\bf V}.$ Обозначим

$$

\Gamma_n({\bf V}) = \Gamma_n/(\Gamma_n \cap \mathop{Id}({\bf V})), \quad \gamma_n({\bf V}) = \mathop{\rm dim} \Gamma_n({\bf V}).

$$



\smallskip



\begin{theorem}[12 \cite{cherev:bib:RSM14}]

Для многообразия алгебр Лейбница"--~Пуассона ${\bf V}$ над произвольным полем следующие условия эквивалентны\/$:$



\begin{itemize}



\item[\rm (i)] последовательность $\{c_n({\bf V})\}_{n \geq 1}$ ограничена полиномом\/$;$



\item[\rm (ii)] для некоторого $m \geq 2$ в ${\bf V}$ выполнены полилинейные тождества

$$

 \{x_1, \dots ,x_m\} = 0,\quad \{x_1,y_1\}\cdot \dots \cdot \{x_m,y_m\} = 0;

$$



\item[\rm (iii)] найдётся такое число $N,$ что для любого $n > N$ выполнено равенство $\gamma_n({\bf V}) = 0;$



\item[\rm (iv)] найдётся такое число $N$, что для любого $n \geq N$ будет выполнено равенство

$$

 c_n({\bf V}) = 1+\sum_{k=2}^NC_n^k \cdot \gamma_k({\bf V}).

$$



\end{itemize}



\end{theorem}



\smallskip



Рассмотрим двумерную алгебру Лейбница $L_2$ над полем $K$ с базисом $a$, $b$ и таблицей умножения $[a,b]=a$, $[a,a]=[b,b]=[b,a]=0$. Обозначим через $B^2$ алгебру Лейбница"--~Пуассона $L_2 \oplus K$ с операциями \eqref{rps10}.





\smallskip



\begin{theorem}[13 \cite{cherev:bib:RSM14}]

Для алгебры Лейбница"--~Пуассона $B^2$ верны следующие утверждения\/$:$



\begin{itemize}



\item[\rm (i)] полилинейные тождества

$$

\{x_1,x_2\}\cdot\{x_3,x_4\}=0,\quad  \{x_1,\{x_2,x_3\}\}=0

$$

порождают идеал тождеств алгебры $B^2;$



\item[\rm (ii)] рост многообразия $\mathop{var}(B^2),$ порождённого алгеброй $B^2,$ является почти полиномиальным$,$ причём для любого натурального $n$ выполнено равенство

$$

c_n(B^2) = n\cdot 2^{n-1}-n+1;

$$



\item[\rm (iii)] базис полилинейной компоненты $P_n(B^2)$ состоит из элементов вида

$$

x_{i_1}\cdot \dots \cdot x_{i_s}\cdot \{x_{j_1}, \dots ,x_{j_t}\},

$$

где $\{i_1, \dots ,i_s,j_1, \dots ,j_t\}=\{1, \dots ,n\}$, $i_1< \dots <i_s$, $j_2< \dots <j_t.$

\end{itemize}

\end{theorem}



\smallskip



Обозначим через ${\bf W}_s$ многообразие алгебр Лейбница"--~Пуассона, порождённое полилинейным тождеством

$$

\{x_1,y_1\}\cdot \dots \cdot \{x_s,y_s\} = 0.

$$



\smallskip





\begin{theorem}[14 \cite{cherev:bib:RSM18}]

Для многообразия алгебр Лейбница"--~Пуассона ${\bf V}$ следующие условия эквивалентны\/$:$



\begin{itemize}



\item[\rm (i)] последовательность $\{c_n({\bf V})\}_{n\geq 1}$ ограничена полиномом\/$;$



\item[\rm (ii)] $U_2 \notin {\bf V},$ $B^2 \notin {\bf V}$ и для некоторого $s\geq 2$ выполнено включение ${\bf V} \subset {\bf W}_s;$



\item[\rm (iii)] существует такая константа $C,$ что в сумме \eqref{dop10}

 $m_\lambda({\bf V}) = 0$ в случае$,$ если выполнено условие $n -\lambda_1 > C.$

\end{itemize}



\end{theorem}





\begin{thebibliography}{99}





\Bibitem{cherev:bib:G-MSP}

\paper Polynomial growth  of the codimentions: a characterization

\by A.~Giambruno, S.~P.~Mishchenko

\jour Proc. Amer. Math. Soc.

\yr 2010

\vol 138

\issue 3

\pages 853--859

\crossref{http://dx.doi.org/10.1090/S0002-9939-09-10160-0}

\mathscinet{http://www.ams.org/mathscinet-getitem?mr=2566551}

\zmath{http://zbmath.org/?q=an:1200.17003}





\Bibitem{cherev:bib:Dr}

\book Free algebras and PI-algebras. Graduate course in algebra

\by V.~Drensky

\publaddr Singapore

\publ Springer-Verlag

\yr 2000

\totalpages xii+271

\mathscinet{http://www.ams.org/mathscinet-getitem?mr=1712064}

\zmath{http://zbmath.org/?q=an:0936.16001}



\RBibitem{cherev:bib:Kemer}

\paper Шпехтовость T-идеалов со степенным ростом коразмерностей

\by А.~Р.~Кемер

\jour Сиб. матем. журнал

\yr 1978

\vol 19

\issue 1

\pages 54--69

\mathscinet{http://www.ams.org/mathscinet-getitem?mr=466190}

\zmath{http://zbmath.org/?q=an:0385.16009}

\misctransl

\paper Spechtianess of T-ideals with polynomial growth of the co-dimensions

\jour Sib. Mat. Zh. 

\vol 19

\issue 1

\pages 54--69 

\yr 1978

\by A.~R.~Kemer







\RBibitem{cherev:bib:M2}

\paper Многообразия алгебр Ли с двуступенно нильпотентным коммутантом

\by С.~П.~Мищенко

\jour Весці Акадэміі навук БССР. Сер. фiз.-мат. навук

\yr 1987

\issue 6

\pages 39--43

\misctransl

\paper Varieties of Lie algebras with two–step nilpotent commutant

\jour Vestsi Akad. Navuk BSSR. Ser. Fiz.-Mat. Navuk 

\yr 1987

\issue 6

\pages 39–43



\RBibitem{cherev:bib:B-Z}

\paper $T$-идеалы свободных алгебр Ли с полиномиальным ростом последовательности коразмерностей

\by И.~И.~Бенедиктович, А.~Е.~Залесский

\jour Весці Акадэміі навук БССР. Сер. фiз.-мат. навук.

\yr 1980

\issue 3

\pages 5--10

\mathscinet{http://www.ams.org/mathscinet-getitem?mr=582766}

\zmath{http://zbmath.org/?q=an:0434.17009}

\misctransl

\by I.~I.~Benediktovich, A.~E.~Zalesskiy

\paper $T$-ideals of free Lie algebras with polynomial growth of a sequence of codimensionalities

\jour Vestsi Akad. Navuk BSSR Ser. Fiz.-Mat. Navuk 

\yr 1980

\issue  3

\pages 5–10



\RBibitem{cherev:bib:M5}

\by С.~П.~Мищенко

\paper О~многообразиях полиномиального роста алгебр~Ли над полем характеристики нуль

\jour Матем. заметки

\yr 1986

\vol 40

\issue 6

\pages 713--721

\mathnet{http://mi.mathnet.ru/mz5228}

\mathscinet{http://www.ams.org/mathscinet-getitem?mr=887447}

\zmath{http://zbmath.org/?q=an:0624.17007}

\transl англ. пер.:~

\paper Varieties of polynomial growth of Lie algebras over a field of characteristic zero

\by S.~P.~Mishchenko 

\jour Math. Notes

\yr 1986

\vol 40

\issue 6

\pages 901--905

\crossref{http://dx.doi.org/10.1007/BF01158347}

\isi{http://gateway.isiknowledge.com/gateway/Gateway.cgi?GWVersion=2&SrcApp=PARTNER_APP&SrcAuth=LinksAMR&DestLinkType=FullRecord&DestApp=ALL_WOS&KeyUT=A1986J487800017}



\RBibitem{cherev:bib:M-V}

\paper A Leibniz variety with almost polynomial growth

\by S.~Mishchenko, A.~Valenti

\jour J. Pure Appl. Algebra

\yr 2005

\vol 202

\issue 1--3

\pages 82--101

\crossref{http://dx.doi.org/10.1016/j.jpaa.2005.01.013}

\mathscinet{http://www.ams.org/mathscinet-getitem?mr=2163402}

\zmath{http://zbmath.org/?q=an:1179.17002}

%\elib{http://elibrary.ru/item.asp?id=13495431}





\RBibitem{cherev:bib:M-Ch}

\by С.~П.~Мищенко, О.~И.~Череватенко

\paper Необходимые и достаточные условия полиномиальности роста многообразия алгебр Лейбница

\jour Фундамент. и прикл. матем.

\yr 2006

\vol 12

\issue 8

\pages 207--215

\mathnet{http://mi.mathnet.ru/fpm37}

\mathscinet{http://www.ams.org/mathscinet-getitem?mr=2314032}

\zmath{http://zbmath.org/?q=an:1184.17003}

%\elib{http://elibrary.ru/item.asp?id=11143844}

\transl англ. пер.:~

\paper Necessary and sufficient conditions for a variety of Leibniz algebras to have polynomial growth

\by S.~P.~Mishchenko, O.~I.~Cherevatenko 

\jour J.~Math. Sci.

\yr 2008

\vol 152

\issue 2

\pages 282--287

\crossref{http://dx.doi.org/10.1007/s10958-008-9054-y}

%\elib{http://elibrary.ru/item.asp?id=13574010}



\RBibitem{cherev:bib:M-Ch2}

\paper Многообразия алгебр Лейбница слабого роста

\by С.~П.~Мищенко, О.~И.~Череватенко

\yr 2006

\pages \mbox{19--23}

\jour Вестн. СамГУ. Естественнонаучн. сер.

\yr 2012

\issue 9(49)

%\elib{http://elibrary.ru/item.asp?id=9505083}

\misctransl

\paper Variety of Leibniz algebras of weak growth 

\by S.~P.~Mishchenko, O.~I.~Cherevatenko

\yr 2006

\pages \mbox{19--23}

\jour Vestnik SamGU. Estestvenno-Nauchnaya Ser.

\yr 2012

\issue 9(49)





\RBibitem{cherev:bib:ChOI}

\paper О нильпотентных алгебрах Лейбница

\by О.~И.~Череватенко

\jour Научные ведомости БелГУ. Математика. Физика

\yr 2012

\issue 17(136)

\pages 132--136

\misctransl

\paper On nilpotent Leibnitz algebras

\by O.~I.~Cherevatenko

\jour Nauchnyye vedomosti BelGU. Matematika. Fizika

\yr 2012

\issue 17(136)

\pages 132--136









\RBibitem{cherev:bib:MPR}

\paper Poisson PI algebras

\by S.~P.~Mishchenko, V.~M.~Petrogradsky, A.~Regev

\jour Trans. Amer. Math. Soc.

\yr 2007

\vol 359

\issue 10

\pages 4669--4694

\crossref{http://dx.doi.org/10.1090/S0002-9947-07-04008-1}

\mathscinet{http://www.ams.org/mathscinet-getitem?mr=2320646}

\zmath{http://zbmath.org/?q=an:1156.17005}

%\elib{http://elibrary.ru/item.asp?id=15292625}





\RBibitem{cherev:bib:RSM12}

\paper Эквивалентные условия полиномиальности роста многообразий алгебр Пуассона

\by С.~М.~Рацеев

\jour Вестн. Моск. унив. Сер. 1. Математика. Механика

\yr 2012

\vol 67

\issue 5

\pages 8--13

\mathscinet{http://www.ams.org/mathscinet-getitem?mr=3076492}

\zmath{http://zbmath.org/?q=an:06190764}

%\elib{http://elibrary.ru/item.asp?id=18260421}

\transl англ. пер.:~

\mathscinet{http://www.ams.org/mathscinet-getitem?mr=3076492}

\zmath{http://zbmath.org/?q=an:06190764}

%\elib{http://elibrary.ru/item.asp?id=20487206}

\by S.~M.~Ratseev

\paper Equivalent conditions of polynomial growth of a variety of Poisson algebras

\jour Mosc. Univ. Math. Bull. 

\vol 67

\issue 5--6

\pages 195--199 

\yr 2012



\RBibitem{cherev:bib:RSM13}

\by С.~М.~Рацеев

\paper Алгебры Пуассона полиномиального роста

\jour Сиб. матем. журн.

\yr 2013

\vol 54

\issue 3

\pages 700--711

\mathnet{http://mi.mathnet.ru/smj2452}

%\elib{http://elibrary.ru/item.asp?id=19080912}

\mathscinet{http://www.ams.org/mathscinet-getitem?mr=3112625}

\transl англ. пер.:~

\paper Poisson algebras of polynomial growth

\by S.~M.~Ratseev 

\jour Siberian Math.~J.

\yr 2013

\vol 54

\issue 3

\pages 555--565

\crossref{http://dx.doi.org/10.1134/S0037446613030191}

%\elib{http://elibrary.ru/item.asp?id=20449177}

\isi{http://gateway.isiknowledge.com/gateway/Gateway.cgi?GWVersion=2&SrcApp=PARTNER_APP&SrcAuth=LinksAMR&DestLinkType=FullRecord&DestApp=ALL_WOS&KeyUT=000322243600019}



\RBibitem{cherev:bib:RSM14}

\by С.~М.~Рацеев

\paper Коммутативные алгебры Лейбница--Пуассона полиномиального роста

\jour Вестн. СамГУ. Естественнонаучн. сер.

\yr 2012

\issue 3/1(94)

\pages 54--65

\mathnet{http://mi.mathnet.ru/vsgu5}

%\elib{http://elibrary.ru/item.asp?id=18242261}

\misctransl

\by S.~M.~Ratseev

\paper Commutative Leibniz--Poisson algebras of polynomial growth

\jour Vestnik SamGU. Estestvenno-Nauchnaya Ser.

\yr 2012

\issue 3/1(94)

\pages 54--65





\RBibitem{cherev:bib:RSM18}

\paper Необходимые и достаточные условия полиномиальности роста многообразий алгебр Лейбница"--~Пуассона

\by С.~М.~Рацеев

\jour Изв. вузов. Матем.

\yr 2014

\issue 3

\pages 33–39

\transl англ. пер.:~

\jour Russian Math. (Iz. VUZ)

\yr 2014

\vol 58

\issue 3

\by S.~M.~Ratseev

\paper Necessary and sufficient conditions of polynom ial growth of varieties of Leibniz--Poisson algebras

\toappear





\end{thebibliography}







\noindent

\hskip6.5cm \thedate \par\noindent

\hskip6.5cm\therevdate



\vspace{-10mm}



%\chead[\fancyplain{}{\scriptsize \sl C\,h\,e\,r\,e\,v\,a\,t\,e\,n\,k\,o~O.\,I.}]{\fancyplain{}{\scriptsize \sl  Varieties of linear algebras of polynomial growth}}



\makeengtitle



\renewcommand{\baselinestretch}{0.9}





\newpage



%\end{document}

